\documentclass{beamer}
%[aspectratio=169]   \usepackage[czech]{babel}
\usepackage{apo-lecture-cz}
\usepackage{pdfpages}
\usepackage{pdfcomment}
\usepackage{listings}
\usepackage{array,multirow}
\usepackage{hhline}

\subtitle{Lekce 01. Úvod}
\author{Pavel Píša \phantom{xxxxxxxxx} Petr Štěpán \\ \small\texttt{pisa@fel.cvut.cz}\phantom{xxxx}\small\texttt{stepan@fel.cvut.cz}}

\begin{document}

\maketitle

\section{Úvod}


\begin{frame}
\frametitle{Motivace}

Co můžete udělat pro zrychlení Vašeho programu:
\begin{itemize}
\item Použít výkonnější počítač:
  \begin{itemize}
  \item Zvýšit výkon CPU 
    \begin{itemize}
    \item Frekvence CPU
    \item Efektivita CPU - kolik operací stihne provést za 1 takt procesoru
    \end{itemize}
  \end{itemize}
\item Změnit program:
  \begin{itemize}
  \item Zlepšit efektivitu práce s pamětí
  \item Paralelizovat program:
    \begin{itemize}
    \item Zvýšit počet jader CPU
    \end{itemize}
  \end{itemize}
\end{itemize}
\end{frame}


\begin{frame}
\frametitle{Motivace}
Má paralelizace nějaká omezení:
\begin{itemize}
\item Nikdy nelze paralelizovat celý program
\item Amdahlův zákon
  \begin{itemize}
  \item $\alpha$ procent programu nelze paralelizovat
  \item zbytek programu $1-\alpha$ lze na $p$ procesorech zrychlit na $\frac{1 - \alpha}{p}$
  \item Poměr zrychlení na $p$ procesorech $S(p) = \frac{\alpha + 1 - \alpha}{\alpha + \frac{1-\alpha}{p}} = \frac{p}{1+\alpha \cdot (p-1)}$
  \item Limit zrychlení je pro nekonečně procesorů $S(p\to\infty) = lim_{p\to\infty}\frac{p}{1+\alpha \cdot (p-1)} = \frac{1}{\alpha}$
    \begin{itemize}
    \item Např pro $\alpha=0.3$ je limit $S(p\to\infty) = 3.\overline{3}$, tedy program nelze zrychlit více než $3.\overline{3}$ krát. 
    \end{itemize}
  \end{itemize}
\item V praxi, čím více máte procesů, tím stoupá i náročnost přípravy dat pro paralelizaci.
\end{itemize}
\end{frame}


\begin{frame}
\frametitle{Motivace}
Proč studovat architektury počítačů:
\begin{itemize}
\item Naučit se jak pracuje počítač, který vykonává Váš program a kde jsou možnosti program zefektivnit 
  \begin{itemize}
  \item zjistit, v čem současný počítač zpomaluje výpočet (rychlost procesoru, velikost vyrovnávací paměti, latence hlavní paměti, počet výpočetních jader) a otestovat jiný HW
  \item zjistit, zda lze program upravit, aby využíval lépe dostupné prostředky 
  \begin{itemize}
    \item upravit pořadí přístupu do paměti a tím lépe využívat vyrovnávací paměť
    \item upravit program, aby požíval méně skoků a tím se vykonával rychleji
    \item paralelizovat výpočet, využít specializovaný HW - GPU, externí výpočetní jednotku např. Coral USB nebo Intel Neural Compute Stick~2.
  \end{itemize}
\end{itemize}
\item Poptávka po absolventech kombinující umělou inteligenci a vestavné systémy (embedded system)
\item Pokud je pro Vás počítač BlackBox, pak jsou Vaše programy neefektivní
\end{itemize}
\end{frame}


\begin{frame}
\frametitle{Obsah přednášek}
Projdeme si všechny základní součásti počítače:
\begin{itemize}
\item CPU
\item Hierarchie paměti - Cache/RAM/Disk
\item Vstupy a výstupy - I/O klávesnice, myš, obrazovka, síťová karta, ovladače HW
\item Výjimky a přerušení - efektivní spolupráce CPU s HW
\end{itemize}

Motivace proč navštěvovat přednášky:
\begin{itemize}
\item Něco se naučíte a budete to mít lehčí při přípravě na zkoušku
\item Kdo správně zodpoví kvízovou otázku v závěru přednášky, dostane bod aktivity
  \begin{itemize}
  \item Kvízy aktivity na přednáškách budou náhodně, první bodovaný kvíz příští týden
  \item Celkem 8 bodů za aktivitu na přednáškách
  \end{itemize}
\end{itemize}

\end{frame}

\begin{frame}
\frametitle{Náplň cvičení}

\begin{columns}
\begin{column}{0.6\textwidth}
\begin{itemize}
\item 4 menší domácí úkoly - 36 bodů
\begin{itemize}
\item 2 programy v C
\item 2 formuláře
\item alespoň 3 úlohy ze 4
\end{itemize}
\item Semestrální úloha - 24 bodů
\begin{itemize}
\item Týmový projekt - dvojice, nebo jednotlivci
\item Speciální HW MZ APO deska
\end{itemize}
\end{itemize}
\end{column}
\begin{column}{0.35\textwidth}  
   \begin{tabular}{|l|l|}\hline
   Známka & Body\\ \hline
   A & >=90\\ \hline
   B & 80 -- 89.9\\ \hline
   C & 70 -- 79.9\\ \hline
   D & 60 -- 69.9\\ \hline
   E & 50 -- 59.9\\ \hline
   F & <50\\ \hline
   \end{tabular}
\end{column}
\end{columns}
\begin{itemize}
\item Nepovinné úlohy nebo aktivita při cvičení/přednáškách - 10 bodů
\end{itemize}
\bigskip
Zkouška: 
\begin{itemize}
\item písemný test 30 bodů, min 15 bodů
\item ústní $\pm$ 10 bodů    
\end{itemize}

\end{frame}


\begin{frame}
\frametitle{Navazující předměty}
Pokud Vás tento předmět zaujme, tak na něj navazují tyto předměty:
\begin{itemize}
\item B4M35PAP - Pokročilé architektury počítačů
\item B3B38VSY - Vestavné systémy
\item B4M38AVS - Aplikace vestavných systémů
\item B4B35OSY - Operační systémy
\item B0B35LSP - Logické systémy a procesory
\end{itemize}
\end{frame}


\begin{frame}
\frametitle{Materiály k předmětu}
\begin{itemize}
\item PATTERSON, David A. a John L. HENNESSY. Computer organization and design RISC-V edition:
    the hardware/software interface. Second Edition. Cambridge: Elsevier, [2021].
    ISBN 978-0-12-820331-6. (12 kusů v ústřední knihovně ČVUT)
\item web:
\begin{itemize}
\item https://cw.fel.cvut.cz/b192/courses/b35apo/
\item https://dcenet.felk.cvut.cz/apo/
\end{itemize}
\item Kurzy v angličtině:
\begin{itemize}
\item MIT 6.004/6.191 – Computation Structures
\item Computation Structures | Electrical Engineering and Computer Science | MIT OpenCourseWare (2015)
\item Computer System Architecture | Electrical Engineering and Computer Science | MIT OpenCourseWare (2005)
\end{itemize}
\item Kurzy v češtině:
\begin{itemize}
\item https://courses.fit.cvut.cz/BI-APS/
\item https://www.vut.cz/studenti/predmety/detail/218515?apid=218515
\end{itemize}

\end{itemize}
\end{frame}

\section{Složení počítače}
\begin{frame}
\frametitle{Co je uvnitř počítače}

Základní deska počítače:
\begin{center}
   \includegraphics[width=0.8\textwidth]{common/computer/motherboard.jpg}
\end{center}

\end{frame}

\begin{frame}
\frametitle{Co je uvnitř počítače}

Rozebraný telefon:
\begin{center}
   \includegraphics[width=0.6\textwidth]{common/computer/mobile.jpg}
\end{center}
\begin{center}
   \includegraphics[width=0.4\textwidth]{common/computer/mobile-cpu.jpg}
\end{center}
\end{frame}


\begin{frame}  
\frametitle{von Neumannova architektura}

Společný koncept navržený maďarským fyzikem Johnem von Neumannem (1903-1957) obsahuje:
\begin{itemize}
\item Procesor - Central Processing Unit - CPU
\item Paměť - Memory, Random-access Memory, 
\item Vstup/Výstup - Input/Output
\end{itemize}
\begin{center}
   \includegraphics[width=0.7\textwidth]{generated/computer/cpu-vonNeumann.pdf}
\end{center}

Základem je jdna paměť použitá jak pro program, tak pro data.

\end{frame}


\begin{frame}
\frametitle{Procesor - CPU}
\begin{columns}
\begin{column}{0.5\textwidth}
\begin{itemize}
\item CPU obsahuje uživatelské registry, které obsahují 16, 32 nebo 64 bitové hodnoty podle architektury čipu
\item CPU načítá z paměti instrukce tak jak jdou za sebou a každou načtenou instrukci vykoná
\item Adresa právě vykonávané instrukce je ve speciálním registru PC nebo IP
\item ALU část CPU která umí sčítat, odčítat, násobit, dělit a další aritmetické operace
\end{itemize}
\end{column}
\begin{column}{0.5\textwidth}  
\begin{center}
   \includegraphics[width=0.95\textwidth]{generated/computer/cpu-cz.pdf}
\end{center}
\end{column}
\end{columns}

\end{frame}


\begin{frame}[fragile]
\frametitle{Paměť}

Paměť si pamatuje data - bajty, slova.

Pokud už znáte nějaký programovací jazyk, tak si ji můžete představit jako pole, např. v jazyce C jako:\\
\begin{minted}{c}
unsigned char RAM[16 * 1024 * 1024 *1024]; // 16GiB RAM
\end{minted}

Z paměti lze číst, například zde do 32-bitového registru R10:\\
\begin{minted}{c}
registr R10 = RAM[adresa] | (RAM[adresa+1]<<8)
           | (RAM[adresa+2]<<16) | (RAM[adresa+3]<<24);
\end{minted}
\bigskip

nebo zapsat registr R10 do paměti:\\
\begin{minted}{c}
RAM[adresa] = R10 & 0xFF; RAM[adresa+1] = (R10>>8) & 0xFF;
RAM[adresa+2] = (R10>>16) & 0xFF;
RAM[adresa+3] = (R10>>24) & 0xFF;
\end{minted}

\end{frame}


\begin{frame}[fragile]
\frametitle{Paměť}

Big endian čtení z paměti probíhá následujícím způsobem:\\
\begin{minted}{c}
registr R10 = RAM[adresa+3] | (RAM[adresa+2]<<8)
           | (RAM[adresa+1]<<16) | (RAM[adresa]<<24);
\end{minted}
\bigskip

nebo zápis do paměti takto:\\
\begin{minted}{c}
RAM[adresa+3] = R10 & 0xFF; RAM[adresa+2] = (R10>>8) & 0xFF;
RAM[adresa+1] = (R10>>16) & 0xFF;
RAM[adresa] = (R10>>24) & 0xFF;
\end{minted}

Některé starší procesory umožňovali pouze čtení ze zarovnané adresy, tedy z adresy adresa=4*x.

Pokud se bude číst z adresy, která není násobkem 4 pak se přečtou dvě slova a z nich se sestaví výsledek.

\end{frame}


\begin{frame}
\frametitle{Instrukce procesoru}

\begin{itemize}
\item Instrukce procesoru jsou jediné činnosti, které procesor umí vykonat
\item Základní instrukce jsou:
\begin{itemize}
\item uložit konstantu do registru
\item načíst data z paměti do registru
\item provést matematickou operaci s registry a výsledek uložit do registru
\item uložit data z registru do paměti
\item porovnat hodnotu dvou registrů
\item podle výsledku porovnání provádět jiné instrukce (změnit PC na jinou hodnotu než je následující instrukce = provést skok v programu)
\end{itemize}
\item veškeré programy, ať už vytvořené v jazyku C, Python, programy provádějící i velmi komplexní výpočty jsou složené pouze z těchto jednoduchých instrukcí procesoru
\end{itemize}

\end{frame}

\begin{frame}
\frametitle{Instrukce procesoru}

Instruction Set Architecture (ISA) Architektura souboru instrukcí
\begin{itemize}
\item je kompletní instrukční sada včetně adresních módů
\item např. x86 (IA-32), x86-64 (AMD64, EM64T, IA-32e), ARM32, ARM64, AVR, MIPS, RISC-V
\item ISA obsahuje:
\begin{itemize}
\item množinu strojových instrukcí procesoru
\item podporované datové typy (celá čísla, celá čísla se znaménkem, reálná čísla)
\item množinu registrů
\item adresní módy
\item organizace paměti
\end{itemize}
\end{itemize}

\end{frame}


\begin{frame}
\frametitle{Instrukce procesoru}

Dva základní přístupy k návrhu ISA
\begin{columns}
\begin{column}{0.5\textwidth}
\begin{center}
\LARGE{RISC}
\end{center}
Reduced Instruction Set Computer
\begin{itemize}
\item Obvykle menší počet instrukcí
\item Všechny instrukce mají stejný počet bajtů, některé umožňují poloviční délku
\item Méně jednodušších adresních módů
\item Matematické operace ALU pouze s registry
\end{itemize}
\end{column}
\begin{column}{0.5\textwidth}  
\begin{center}
\LARGE{CISC}
\end{center}
Complex Instruction Set Computer
\begin{itemize}
\item Obvykle větší počet instrukcí
\item Délka instrukcí je od 1 bajtu do např. 12 bajtů, nejčastější instrukce jsou nejkratší
\item Obvykle mnoho i složitých adresních módu
\item Matematické oprace ALU i s hodnotami z paměti
\end{itemize}

\end{column}
\end{columns}
\end{frame}


\begin{frame}
\frametitle{Překlad programu}

Jak je možné, že jste o strojových instrukcích ještě neslyšeli?

\begin{center}
   \includegraphics[width=0.65\textwidth]{generated/overview/preklad.pdf}
\end{center}

\begin{itemize}
\item Programování v assembleru (jazyku symbolických adres) je nekomfortní.
\item Překladač přeloží vyšší programovací jazyk přímo do strojových instrukcí procesoru
\item Křízový překlad je, když přeložíte program pro jiný typ procesoru, než na kterém probíhá překlad.
\end{itemize}
\end{frame}



\begin{frame}
\frametitle{Procesor}

Z čeho se vlastně procesor skládá?
\begin{itemize}
\item Možná jste někdy slyšeli, že procesor obsahuje třeba 16 miliard tranzistorů
\item V roce 1965 spoluzakladatel společnosti Intel Gordon Moor formuloval zákon: 
\begin{itemize}
\item Počet tranzistorů, které mohou být umístěny na integrovaný obvod se při zachování stejné ceny zhruba každých 18 měsíců zdvojnásobí.
\end{itemize}
\item Více méně platí až do dnes, i když už se dostáváme na hranice fyzikálních možností.
\end{itemize}

K čemu potřebujeme tranzistory v procesoru (CPU)?
\begin{itemize}
\item Abychom implementovali Booleovu algebru.
\end{itemize}
\end{frame}



\section{Logické obvody - Opakování PSY}

\begin{frame}
\frametitle{Binární komparátor}

Jak udělat komparátor na rovnost?
\begin{itemize}
\item Máme dvě 64-bitová čísla X a Y která potřebujeme porovnat
\item Porovnání dvou bitů provede operace xnor, která má hodnotu 1 pokud dva bity mají stejnou hodnotu.
\begin{itemize}
\item Nyní potřebujeme udělat and 64 logických hodnot.
\end{itemize}
\item Jak to můžeme udělat?
\end{itemize}
\end{frame}


\begin{frame}
\frametitle{Lineární řešení}

Řešení může být lineární:
\begin{center}
   \includegraphics[width=0.65\textwidth]{generated/logic/lin-comparator.pdf}
\end{center}

Jaké je to řešení?
\begin{itemize}
\item Počet and hradel se dvěma vstupy - 63 hradel
\item Rychlost výpočtu?
\begin{itemize}
\item Rychlost výpočtu hradla 4ps, protože na sebe výpočty navazují potřebujeme 63 výpočtů
\item Rychlost výpočtu komparátoru bude 252 ps, což odpovídá frekvenci 4GHz  to je velmi polmalé
\end{itemize}
\end{itemize}
\end{frame}


\begin{frame}
\frametitle{Logaritmické řešení}

\begin{columns}
\begin{column}{0.48\textwidth}
Rychlejší řešení využívá paralelního výpočtu všech dvojic a následné spojování do výsledku:
\begin{center}
   \includegraphics[width=0.95\textwidth]{generated/logic/log-comparator.pdf}
\end{center}
\end{column}
\begin{column}{0.52\textwidth}

Jaké je to řešení?
\begin{itemize}
\item Počet and hradel se dvěma vstupy - 63 hradel
\begin{itemize}
\item 32 + 16 + 8 + 4 + 2 + 1 = 63
\end{itemize}
\item Rychlost výpočtu?
\begin{itemize}
\item Výpočty v jedné řadě jsou nezávislé a probíhají ve stejnou dobu
\item Kolik řad máme? První řada obsahuje 32 hradel and, druhá 16 and, třetí 8 and, čtvrtá 4 and, pátá 2 and, šestá 1 and
\item Rychlost výpočtu komparátoru bude 6*4 = 24 ps, což je 10 krát rychlejší
\end{itemize}
\end{itemize}
\end{column}
\end{columns}

\end{frame}


\begin{frame}[fragile]
\frametitle{Rychlé mocnění}
Výpočet mocnin - opakování z PSY:
\begin{columns}
\begin{column}{0.44\textwidth}
\small
\begin{minted}{c}
double Exp(double a, int k) {
  if (k == 0) return 1;
  return a * Exp(a, k - 1);
}
\end{minted}
\bigskip

Lineární mocnění - v kryptografii je potřeba vypočítat i výrazy $X^{1000000}$

K tomu je potřeba provést milión násobení.

\end{column}
\hfill
\begin{column}{0.50\textwidth}
\small
\begin{minted}{c}
double FastExp(double a, int k) {
  if (k == 0) return 1;
  if (k % 2 == 0) {
    double i = FastExp(a, k / 2);
    return i * i;
  } else {
    return a * FastExp(a, k - 1);
  }
}
\end{minted}
\bigskip

Logaritmické mocnění - nahradí milión násobení, desítkami násobení.

Logaritmické mocnění je 100000 krát rychlejší.

\end{column}
\end{columns}
\end{frame}




\begin{frame}
\frametitle{Sčítání - Ripple Carry Adder - opakování}


Lineární sčítačka se jmenuje Ripple Carry Adder (Řetězová sčítačka).

\begin{center}
   \includegraphics[width=0.7\textwidth]{generated/logic/addition/ripple_adder.pdf}
\end{center}

Je obdobně pomalá jako lineární komparátor - vlastně 2-krát tak pomalá než lineární komparátor

\end{frame}


\begin{frame}
\frametitle{Sčítání - Carry Lookahead Adder}
Přímý výpočet přenosů -- carry na základě hodnot bitů sčítanců, je nejrychlejší sčítačka pro 4-bitová čísla (opakování PSY).

Dvě čísla sečte na 4 zpoždění hradla - teoreticky konstantní doba vzhledem k počtu bitů čísel.
\begin{center}
   \includegraphics[width=0.5\textwidth]{generated/logic/addition/lookahead_direct.pdf}
\end{center}

Pro 64 bitová čísla by potřebovala sčítačka kolem $10^{20}$ tranzistorů, což je moc.
\end{frame}



\begin{frame}
\frametitle{Sčítání - Carry Lookahead Adder}

Co tedy s většími čísly, třeba 64-bitovými?

\begin{columns}
\begin{column}{0.45\textwidth}
\begin{itemize}
\item Upravíme 4-bitovou sčítačku tak, aby mohla přijmout přenos $C$ z předchozích výpočtů.
\begin{itemize}
\item POZOR budou se muset přidat hradla pro zpracování carry od sčítaček nižších řádů.
\end{itemize}
\item Takto upravené sčítačky můžeme řetězit.
\end{itemize}
\end{column}
\begin{column}{0.55\textwidth}
\begin{center}
   \includegraphics[width=0.95\textwidth]{generated/logic/addition/lookahead_ripple.pdf}
\end{center}
\end{column}
\end{columns}

\begin{itemize}
\item Zrychlení je pouze o konstantu, závislost na počtu bitů sčítanců je pořád lineární.
\item Rychlost 16-ti bitové sčítačky z obrázku bude 16 zpoždění hradel
\item Rychlost 64-ti bitové sčítačky bude 64 zpoždění hradel, což je 2 krát lepší, než 127 zpoždění u jednoduché Ripple carry adder.
\end{itemize}
\end{frame}


\begin{frame}
\frametitle{Sčítání - Carry Select Adder}
Jiné řešení?

\begin{itemize}
\item Víme, že Carry je buď 0 nebo 1
\item 4bit CLA zdvojíme a spočteme výsledek pokud vstupní Carry je 0 i 1
\item Nakonec řetězíme pouze multiplex, zda skutečně Carry bylo 0 nebo 1
\end{itemize}
\begin{center}
   \includegraphics[width=0.75\textwidth]{generated/logic/addition/lookahead_select.pdf}
\end{center}

\begin{itemize}
\item Sčítačka je rychlejší opět o konstantu krát, místo zpoždění 4 hradel bude řetězit pouze zpoždění 2 hradel pro multiplexor
\item Rychlost 16-ti bitové sčítačky z obrázku bude 10 zpoždění hradel
\item Rychlost 64-ti bitové sčítačky bude 34 zpoždění hradel, což je 2 krát lepší, než 64 zpoždění u řetězení CLA.
\end{itemize}

\end{frame}



\begin{frame}
\frametitle{Propagace a generování carry}

V PSY jsme si ukazovali, že pokud se zaměříme na jeden sloupeček výpočtu součtu, mohou nastat pouze 3 rozdílné situace:
\begin{columns}
\begin{column}{0.25\textwidth}
\begin{center}
   \includegraphics[width=0.45\textwidth]{generated/logic/addition/prop_gen_carry.pdf}
\end{center}
\end{column}
\begin{column}{0.75\textwidth}

\begin{itemize}
\item Pokud je $X_i=0$ a $Y_i=0$, pak musí být $C_{i+1}=0$
\item Pokud je $X_i=1$ a $Y_i=0$, nebo $X_i=0$ a $Y_i=1$, pak musí být $C_{i+1}=C_i$.
\begin{itemize}
\item Tento ppřípad naveme propagování carry, protože $C_{i+1}$ propaguje hodnotu $C_i$
\item $P_i = X_i \oplus Y_i$
\end{itemize}
\item Pokud je $X_i=1$ a $Y_i=1$, pak musí být $C_{i+1}=1$.
\begin{itemize}
\item Tento ppřípad naveme generování carry, protože $C_{i+1}$ nezávisí na $C_i$ a má vždy hodnotu $1$
\item $G_i = X_i \land Y_i$
\end{itemize}
\end{itemize}
\end{column}
\end{columns}

\end{frame}


\begin{frame}
\frametitle{Propagace a generování carry}

Nyní si zkusíme vypočítat hodnotu propagace a generování carry pro skupinu bitů od j do i:
\begin{columns}
\begin{column}{0.45\textwidth}
\begin{center}
   \includegraphics[width=\textwidth]{generated/logic/addition/prop_gen_block-cz.pdf}
\end{center}
\end{column}
\begin{column}{0.55\textwidth}

\begin{itemize}
\item Propagování carry nastane, pokud kombinace hodnot $X_j ... X_i$ a $Y_j ... Y_i$ způsobí, že $C_{i+1}=C_j$.
\begin{itemize}
\item Jak toto $P_{j,i}$ spočítat si ukážeme za chvíli
\end{itemize}
\item Generování carry nastane, pokud kombinace hodnot $X_j ... X_i$ a $Y_j ... Y_i$ způsobí, že $C_{i+1}=1$.
\begin{itemize}
\item Tento případ si označíme $G_{j,i}$
\end{itemize}
\end{itemize}
\end{column}
\end{columns}

\end{frame}


\begin{frame}
\frametitle{Propagace a generování carry}

Jak tedy $P_{j,i}$ a $G_{j,i}$ spočítat? Logaritmicky - rozděl a panuj.

\begin{columns}
\begin{column}{0.5\textwidth}
\begin{center}
   \includegraphics[width=0.95\textwidth]{generated/logic/addition/prop_gen_compute.pdf}
\end{center}
\end{column}
\begin{column}{0.5\textwidth}

\begin{itemize}
\item Hodnoty $G_{i,i}=G_i=X_i \land Y_i$ a $P_{i,i}=P_i=X_i \oplus Y_i$
\item Nyní si musíme uvědomit, kdy dojde k propagaci carry?
\begin{itemize}
\item Pokud dojde k propagaci carry mezi bity j a k a zároveň dojde k propagaci carry mezi bity k+1 a i
\item $P_{j,i}=P_{j,k} \land P_{k+1,i}$ tedy $P_{j,i}=P_h \land P_l$
\end{itemize}
\end{itemize}
\end{column}
\end{columns}

\begin{itemize}
\item Kdy dojde ke generování carry?
\begin{itemize}
\item Pokud vyšší bity vygenerují carry nebo nižší bity vygenerují carry a vyšší bity ho budou propagovat
\item $G_{j,i}=G_{k+1,i} \lor (G_{j,k} \land P_{k+1,i})$ tedy $G_{j,i}=G_h \lor (G_l \land P_h)$
\end{itemize}
\end{itemize}

\end{frame}


\begin{frame}
\frametitle{Propagace a generování carry}

Plán výpočtu hodnot propagace a generování carry.

\begin{center}
   \includegraphics[width=0.45\textwidth]{generated/logic/addition/prop_gen_example.pdf}
\end{center}

\begin{itemize}
\item V prvním kroku spočteme základní hodnoty propagace a generování pro dvojice $X_i$, $Y_i$
\item V dalšm kroku spočteme hodnoty propagace a generování pro čtveřice $X_{2*i},X_{2*i+1}$, $Y_{2*i}, Y_{2*i+1}$
\item V každém dalším kroku se vždy spojí dva bloky do jednoho
\end{itemize}

\end{frame}


\begin{frame}
\frametitle{Sčítání - Carry Lookahead Adder}

Takto lze realizovat výše uvedený výpočet:
\begin{columns}
\begin{column}{0.55\textwidth}
\begin{center}
   \includegraphics[width=0.95\textwidth]{generated/logic/addition/lookahead_tree.pdf}
\end{center}
\end{column}
\begin{column}{0.45\textwidth}

Doba výpočtu dvojic $G_{i,k}$ a $P_{i,k}$:
\begin{itemize}
\item Rychlost 4 bitových sčítanců 5 zpoždění hradel
\item Rychlost 8 bitových sčítanců 7 zpoždění hradel
\item Rychlost 16 bitových sčítanců 9 zpoždění hradel
\item Rychlost 32 bitových sčítanců 11 zpoždění hradel
\item Rychlost 64 bitových sčítanců 13 zpoždění hradel
\end{itemize}

\end{column}
\end{columns}

\end{frame}




\begin{frame}
\frametitle{Sčítání - Carry Lookahead Adder}

Nyní nám zbývá vypočítat všechny Carry $C_i$, pokud bychom předpokládali řetězení, tak známe $C_0$ jinak by :
\begin{itemize}
\item Opět použijeme strom, tentokrát v opačném pořadí než výpočet $P_{i,j}$ a $G_{i,j}$:
\begin{itemize}
\item Pokud známe $C_i$, $G_{j,i}$ a $P_{j,i}$
\item pak $C_{i+1} = G_{j,i} \lor (C_j \land P_{j,i})$
\end{itemize}
\item Výpočet pro 4-bitové sčítance bude probíhat v následujícím pořadí:
\begin{itemize}
\item $C_4 = G_{0,3} \lor (C_0 \land P_{0,3})$, $C_2 = G_{0,1} \lor (C_0 \land P_{0,1})$
\item $C_3 = G_{2,2} \lor (C_2 \land P_{2,2})$, $C_1 = G_{0,0} \lor (C_0 \land P_{0,0})$
\end{itemize}
\item Opět platí, že pro $2\times$ větší sčítance se doba prodlouží o 2 zpoždění. 
\item Tedy pro 64 bitový je výpočet všech $C_i$ s 12 zpožděním hradel.
\item Celý součet pro 64 bitový sčítance trvá 26 zpožděním hradel.
\end{itemize}
\end{frame}



\begin{frame}
\frametitle{Sčítání - Carry Lookahead Adder}

Když máme spočtené G a P pro zadané rozsahy, pak začneme počítat vlastní carry:
\begin{columns}
\begin{column}{0.55\textwidth}
\begin{center}
   \includegraphics[width=0.95\textwidth]{generated/logic/addition/prop_gen_final.pdf}
\end{center}
\end{column}
\begin{column}{0.45\textwidth}

Obecně lze napsat:\\
$C_{i+1} = G_{j,i}$ or ($C_j$ and $P_{j,i}$)

Výpočet C proběhne ve 4 krocích, barevně odlišených:
\begin{itemize}
\item Výpočet $C_8$
\item Výpočet $C_4$
\item Výpočet $C_6$, $C_2$
\item Výpočet $C_7$, $C_5$, $C_3$, $C_1$
\end{itemize}

\end{column}
\end{columns}
\end{frame}


\begin{frame}
\frametitle{Sčítání - Carry Lookahead Adder}

Výpočet $C_i$ a pak i $S_i$ vypadá následovně:
\begin{columns}
\begin{column}{0.5\textwidth}
\begin{center}
   \includegraphics[width=0.95\textwidth]{generated/logic/addition/lookahead_tree_carry.pdf}
\end{center}
\end{column}
\begin{column}{0.5\textwidth}

Doba výpočtu všech carry $C_i$ a $S_i$:
\begin{itemize}
\item Rychlost 4-ti bitových sčítanců 5 zpoždění hradel
\item Rychlost 8-ti bitových sčítanců 7 zpoždění hradel
\item Rychlost 64-ti bitových sčítanců 13 zpoždění hradel
\end{itemize}

Celkem doba výpočtu pro 64-bitové sčítance je 24 zpoždění tedy asi 96ps.

\end{column}
\end{columns}

\end{frame}




\begin{frame}
\frametitle{Sčítání - Carry Lookahead Adder - Příklad}

Výpočet generování a propagace carry bloku pomocí:
\begin{itemize}
\item $g_h$, $p_h$ - generování a propagace carry ve vyšším podbloku
\item $g_l$, $p_l$ - generování a propagace carry v nižším podbloku
\end{itemize}
\begin{table}
\begin{tabular}{|c|c|c|c|c|c|c|c|c|}\hline
a & 0 & 0 & 1 & 0 & 1 & 0 & 0 & 1 \\ \hline
b & 1 & 0 & 1 & 1 & 0 & 1 & 1 & 1 \\ \hhline{|=|=|=|=|=|=|=|=|=|}
g = $a_i$ and $b_i$ & 0 & 0 & 1 & 0 & 0 & 0 & 0 & 1 \\ \hline
p = $a_i$ xor $b_i$ & 1 & 0 & 0 & 1 & 1 & 1 & 1 & 0 \\ \hhline{|=|=|=|=|=|=|=|=|=|}
g = $g_h$ or ( $g_l$ and $p_h$ )& \multicolumn{2}{c|} 0 & \multicolumn{2}{c|} 1 & \multicolumn{2}{c|} 0 & \multicolumn{2}{c|} 1 \\ \hline
p = $p_h$ and $p_l$ & \multicolumn{2}{c|} 0 & \multicolumn{2}{c|} 0 & \multicolumn{2}{c|} 1 & \multicolumn{2}{c|} 0 \\ \hhline{|=|==|==|==|==|}
g = $g_h$ or ( $g_l$ and $p_h$ ) & \multicolumn{4}{c|} 0 & \multicolumn{4}{c|} 1  \\ \hline
p = $p_h$ and $p_l$ & \multicolumn{4}{c|} 0 & \multicolumn{4}{c|} 0  \\ \hhline{|=|====|====|}
g = $g_h$ or ( $g_l$ and $p_h$ ) & \multicolumn{8}{c|} 0   \\ \hline
p = $p_h$ and $p_l$ & \multicolumn{8}{c|} 0   \\ \hline
\end{tabular}
\end{table}
\end{frame}




\begin{frame}{shrink=5}
\frametitle{Sčítání - Carry Lookahead Adder - Příklad}

Výpočet carry na základě carry nižších řádů a příznaků g,p:
\begin{itemize}
\item $C_i = g$ or $(p$ and $C_{i-k})$ - buď se carry vygeneruje a nebo se propaguje z nižších řádů
\end{itemize}
\begin{table}
\footnotesize
\begin{tabular}{|p{0.02\textwidth}|p{0.08\textwidth}|p{0.08\textwidth}|p{0.08\textwidth}|p{0.08\textwidth}|p{0.08\textwidth}|p{0.08\textwidth}|p{0.08\textwidth}|p{0.08\textwidth}|}\hline
a & 0 & 0 & 1 & 0 & 1 & 0 & 0 & 1 \\ \hline
b & 1 & 0 & 1 & 1 & 0 & 1 & 1 & 1 \\ \hhline{|=|=|=|=|=|=|=|=|=|}
g  & 0 & 0 & 1 & 0 & 0 & 0  & 0 & 1 \\ \hline
p & 1 & 0 & 0 & 1   & 1 & 1 & 1 & 0 \\ \hline
C & \multicolumn{2}{p{0.16\textwidth}|}{}  & \multicolumn{2}{p{0.16\textwidth}|}{} & \multicolumn{2}{p{0.16\textwidth}|}{} & \multicolumn{2}{p{0.16\textwidth}|}{}  \\ \hhline{|=|=|=|=|=|=|=|=|=|}
g & \multicolumn{2}{c|} 0 & \multicolumn{2}{c|}{\textcolor{brown}1} & \multicolumn{2}{c|} 0 & \multicolumn{2}{c|} 1 \\ \hline
p  & \multicolumn{2}{c|} 0 & \multicolumn{2}{c|}{\textcolor{teal} 0} & \multicolumn{2}{c|} 1 & \multicolumn{2}{c|} 0 \\ \hline
C  & \multicolumn{2}{l|}{} & \multicolumn{2}{l|}{} & \multicolumn{2}{l|}{} & \multicolumn{2}{l|}{} \\ \hhline{|=|==|==|==|==|}
g  & \multicolumn{4}{c|} 0 & \multicolumn{4}{c|}{\textcolor{cyan}1}  \\ \hline
p  & \multicolumn{4}{c|} 0 & \multicolumn{4}{c|}{\textcolor{green}0}  \\ \hline
C  & \multicolumn{3}{c|}{} & \multicolumn{5}{l|}{$C_4$ = \textcolor{cyan}1 or ( \textcolor{green}0 and $C_0$) = 1}  \\ \hhline{|=|====|====|}
g  & \multicolumn{8}{c|} 0   \\ \hline
p  & \multicolumn{8}{c|} 0   \\ \hline
C  & \multicolumn{8}{l|}{$C_8$ = 0 or ( 0 and $C_0$) = 0}   \\ \hline
\end{tabular}
\end{table}
\end{frame}


\begin{frame}{shrink=5}
\frametitle{Sčítání - Carry Lookahead Adder - Příklad}

Výpočet carry na základě carry nižších řádů a příznaků g,p:
\begin{itemize}
\item $C_i = g$ or $(p$ and $C_{i-k})$ - buď se carry vygeneruje a nebo se propaguje z nižších řádů
\end{itemize}
\begin{table}
\footnotesize
\begin{tabular}{|p{0.02\textwidth}|p{0.08\textwidth}|p{0.08\textwidth}|p{0.08\textwidth}|p{0.08\textwidth}|p{0.08\textwidth}|p{0.08\textwidth}|p{0.08\textwidth}|p{0.08\textwidth}|}\hline
a & 0 & 0 & 1 & 0 & 1 & 0 & 0 & 1 \\ \hline
b & 1 & 0 & 1 & 1 & 0 & 1 & 1 & 1 \\ \hhline{|=|=|=|=|=|=|=|=|=|}
g  & 0 & 0 & 1 & 0 & 0 & 0  & 0 & 1 \\ \hline
p & 1 & 0 & 0 & 1   & 1 & 1 & 1 & 0 \\ \hline
C & \multicolumn{2}{p{0.16\textwidth}|}{}  & \multicolumn{2}{p{0.16\textwidth}|}{} & \multicolumn{2}{p{0.16\textwidth}|}{} & \multicolumn{2}{p{0.16\textwidth}|}{}  \\ \hhline{|=|=|=|=|=|=|=|=|=|}
g & \multicolumn{2}{c|} 0 & \multicolumn{2}{c|}{\textcolor{brown}1} & \multicolumn{2}{c|} 0 & \multicolumn{2}{c|} 1 \\ \hline
p  & \multicolumn{2}{c|} 0 & \multicolumn{2}{c|}{\textcolor{teal} 0} & \multicolumn{2}{c|} 1 & \multicolumn{2}{c|} 0 \\ \hline
C  &  & \multicolumn{3}{l|}{$C_6$ = \textcolor{brown}1 or ( \textcolor{teal}0 and $C_4$) = 1} &  & \multicolumn{3}{l|}{$C_2$ = 1 or ( 0 and $C_0$) = 1} \\ \hhline{|=|==|==|==|==|}
g  & \multicolumn{4}{c|} 0 & \multicolumn{4}{c|}{\textcolor{cyan}1}  \\ \hline
p  & \multicolumn{4}{c|} 0 & \multicolumn{4}{c|}{\textcolor{green}0}  \\ \hline
C  & \multicolumn{3}{c|}{} & \multicolumn{5}{l|}{$C_4$ = \textcolor{cyan}1 or ( \textcolor{green}0 and $C_0$) = 1}  \\ \hhline{|=|====|====|}
g  & \multicolumn{8}{c|} 0   \\ \hline
p  & \multicolumn{8}{c|} 0   \\ \hline
C  & \multicolumn{8}{l|}{$C_8$ = 0 or ( 0 and $C_0$) = 0}   \\ \hline
\end{tabular}
\end{table}
\end{frame}


\begin{frame}{shrink=5}
\frametitle{Sčítání - Carry Lookahead Adder - Příklad}

Výpočet carry na základě carry nižších řádů a příznaků g,p:
\begin{itemize}
\item $C_i = g$ or $(p$ and $C_{i-k})$ - buď se carry vygeneruje a nebo se propaguje z nižších řádů
\end{itemize}
\begin{table}
\footnotesize
\begin{tabular}{|p{0.02\textwidth}|p{0.08\textwidth}|p{0.08\textwidth}|p{0.08\textwidth}|p{0.08\textwidth}|p{0.08\textwidth}|p{0.08\textwidth}|p{0.08\textwidth}|p{0.08\textwidth}|}\hline
a & 0 & 0 & 1 & 0 & 1 & 0 & 0 & 1 \\ \hline
b & 1 & 0 & 1 & 1 & 0 & 1 & 1 & 1 \\ \hhline{|=|=|=|=|=|=|=|=|=|}
g  & 0 & \textcolor{red}0 & 1 & \textcolor{magenta}0 & 0 & \textcolor{brown}0  & 0 & \textcolor{purple}1 \\ \hline
p & 1 & \textcolor{blue}0 & 0 & \textcolor{green}1   & 1 & \textcolor{orange}1 & 1 & \textcolor{teal}0 \\ \hline
C & \multicolumn{2}{p{0.16\textwidth}|}{$C_7$ = \textcolor{red}0 or ( \textcolor{blue}0 and $C_6$) = 0}  & \multicolumn{2}{p{0.16\textwidth}|}{$C_5$ = \textcolor{magenta}0 or ( \textcolor{green}1 and $C_4$) = 1} & \multicolumn{2}{p{0.16\textwidth}|}{$C_3$ = \textcolor{brown}0 or ( \textcolor{orange}1 and $C_2$) = 1} & \multicolumn{2}{p{0.16\textwidth}|}{$C_1$ = \textcolor{purple}1 or ( \textcolor{teal}0 and $C_0$) = 1}  \\ \hhline{|=|=|=|=|=|=|=|=|=|}
g & \multicolumn{2}{c|} 0 & \multicolumn{2}{c|}{\textcolor{brown}1} & \multicolumn{2}{c|} 0 & \multicolumn{2}{c|} 1 \\ \hline
p  & \multicolumn{2}{c|} 0 & \multicolumn{2}{c|}{\textcolor{teal} 0} & \multicolumn{2}{c|} 1 & \multicolumn{2}{c|} 0 \\ \hline
C  &  & \multicolumn{3}{l|}{$C_6$ = \textcolor{brown}1 or ( \textcolor{teal}0 and $C_4$) = 1} &  & \multicolumn{3}{l|}{$C_2$ = 1 or ( 0 and $C_0$) = 1} \\ \hhline{|=|==|==|==|==|}
g  & \multicolumn{4}{c|} 0 & \multicolumn{4}{c|}{\textcolor{cyan}1}  \\ \hline
p  & \multicolumn{4}{c|} 0 & \multicolumn{4}{c|}{\textcolor{green}0}  \\ \hline
C  & \multicolumn{3}{c|}{} & \multicolumn{5}{l|}{$C_4$ = \textcolor{cyan}1 or ( \textcolor{green}0 and $C_0$) = 1}  \\ \hhline{|=|====|====|}
g  & \multicolumn{8}{c|} 0   \\ \hline
p  & \multicolumn{8}{c|} 0   \\ \hline
C  & \multicolumn{8}{l|}{$C_8$ = 0 or ( 0 and $C_0$) = 0}   \\ \hline
\end{tabular}
\end{table}
\end{frame}


\begin{frame}
\frametitle{Odčítání}

Nyní jsme si ukázali optimální sčítačku, pracující i pro čísla o velikost 128 bitů.
\begin{columns}
\begin{column}{0.5\textwidth}
Odčítání lze řešit:
\begin{itemize}
\item speciálním obvodem podobným sčítačce se všemi možnostmi zrychlení jako byly u sčítání
\item nebo z druhého čísla vytvoříme číslo opačné a to pak jednoduše sečteme s tím prvním
\end{itemize}
\end{column}
\hfill
\begin{column}{0.5\textwidth}
\begin{center}
\includegraphics[width=0.5\textwidth]{generated/logic/addition/substaction-cz.pdf}
\end{center}
\end{column}
\end{columns}

V PSY jsme si ukázali obvod pro přičtení 1, jedná se vlastně o výpočet and pro všechny kombinace bitů od nejnižšího bitu.

\end{frame}


\begin{frame}
\frametitle{Násobení celých čísel}

Stejný princip násobení, jak jste se ho naučili na základní škole pro desítkovou soustavu:
\begin{columns}
\begin{column}{0.4\textwidth}
\texttt{\phantom{xxx}153}\\
\texttt{\phantom{xxx}*45}\\
\vspace{-8pt}
\rule[0pt]{1.5cm}{0.1pt}\\
\texttt{\phantom{xxx}765}\\
\texttt{\phantom{xx}612 }\\
\vspace{-8pt}
\rule[0pt]{1.5cm}{0.1pt}\\
\texttt{\phantom{xx}6885}\\
\end{column}
\hfill
\begin{column}{0.4\textwidth}
\texttt{\phantom{xxxxxx}10011001}\\
\texttt{\phantom{xxxxxxx}*101101}\\
\vspace{-8pt}
\rule[0pt]{3cm}{0.4pt}\\
\texttt{\phantom{xxxxxx}10011001}\\
\texttt{\phantom{xxxxx}00000000}\\
\texttt{\phantom{xxxx}10011001}\\
\texttt{\phantom{xxx}10011001}\\
\texttt{\phantom{xx}00000000}\\
\texttt{\phantom{x}10011001}\\
\vspace{-8pt}
\rule[0pt]{3cm}{0.4pt}\\
\texttt{\phantom{x}1101011100101}\\
\end{column}
\end{columns}

\end{frame}

\begin{frame}
\frametitle{Násobení celých čísel}

Podle uvedeného algoritmu můžeme vytvořit následující násobičku s posuvným registrem:\\
(A,B 32 bitů, výsledek 64 bitů)
\begin{center}
\includegraphics[width=0.45\textwidth]{generated/logic/multiplication/multiplier-seq-cz.pdf}
\end{center}
\begin{itemize}
\item Výsledek je ve dvou registrech AC a B
\item Pomalé, už sčítání je náročné, nyní 32 sčítání, tedy přibližně 32*20=640 zpoždění, tedy asi 2560ps, tedy 10 taktů procesoru.
\item Pro násobení 64-bitových čísel přibližně 64*24=1536 zpoždění, tedy asi 25 taktů procesoru.
\end{itemize}
\end{frame}


\begin{frame}
\frametitle{Rychlé násobení - Motivace Wallace tree}

Jak to zrychlit - odložené carry (Carry Save Adder).\\
Carry Save Adder je vlastně Full Adder s tím, že oddělíme součet a přenos do rozdílných čísel.\\
Jak nejrychleji spočítat součet čtyř 32-bitových čísel:\\
\begin{columns}
\begin{column}{0.35\textwidth}
\includegraphics[width=1.0\textwidth]{generated/logic/multiplication/wallace-tree.pdf}
\end{column}
\hfill
\begin{column}{0.65\textwidth}
\begin{itemize}
\item Sečteme paralelně p=w+x a q=y+z a potom p+q -- trvá dlouho, nejméně doby dvou plných součtů
\item Odložíme carry (nebudeme carry propagovat, jen v posledním kroku):
\begin{itemize}
\item 1. krok použijeme Full adder a sečteme vždy bity $w_i+x_i+y_i = c'_ip_i$
\item 2. krok použijeme Full adder a sečteme vždy bity $p_i+c'_{i-1}+z_i = c_iq_i$
\item 3. krok použijeme normální sčítačku 32-bitových čísel ($s_0=q_0$, $c'_{31}$ připojíme ke $q$)
\end{itemize}
\end{itemize}
\end{column}
\end{columns}
\end{frame}


\begin{frame}
\frametitle{Rychlé násobení - Motivace Wallace tree}
Jakého zrychlení tedy dosáhneme?

Sčítání šesti čísel:
\begin{columns}
\begin{column}{0.3\textwidth}
\includegraphics[width=0.98\textwidth]{generated/logic/multiplication/wallace-tree6.pdf}
\end{column}
\hfill
\begin{column}{0.7\textwidth}
\begin{itemize}
\item Každý krok zredukuje velikost problému na $\lceil \frac{2}{3} \rceil$
\item Kolik kroků potřebuje na součet 6 čísel:
\begin{itemize}
\item 1. krok redukce na součet 4 čísel - 2 zpoždění
\item 2. krok redukce na součet 3 čísel - 2 zpoždění
\item 3. krok redukce na součet 2 čísel - 2 zpoždění
\item 4. krok normální sčítačka 2 čísel - 24 zpoždění
\end{itemize}
\item Součet 6 čísel - 30 zpoždění
\item Rychlost součtu N čísel, počet kroků na redukci součtu dvou čísel $\log_{\frac{3}{2}}(N) = \frac{1}{\log(\frac{3}{2})} \cdot \log(N) = $ $ = 5.678 \cdot \log(N) = O(log(N))$
\end{itemize}
\end{column}
\end{columns}
\end{frame}


\begin{frame}
\frametitle{Rychlé násobení - Wallace tree}

Použijeme předchozí princip na co nejrychlejší součet 32 nebo 64 různých čísel:
\begin{columns}
\begin{column}{0.6\textwidth}
\includegraphics[width=1.0\textwidth]{generated/logic/multiplication/wallace-tree2.pdf}
\end{column}
\hfill
\begin{column}{0.4\textwidth}
\begin{itemize}
\item Součiny jsou jednoduché: $x_i \cdot y_j=x_i$~\texttt{and}~$y_j$
\item Nejtěžší je sečíst prostřední sloupec 64 jednobitových čísel
\item Pustíme na všechny bity, co můžeme, paralelně sčítačky a carry budeme přičítat v dalších krocích
\item V první fázi to bude 1323 sčítaček
\end{itemize}
\end{column}
\end{columns}

\end{frame}

\begin{frame}[shrink=5]
\frametitle{Rychlé násobení - Wallace tree}

Když se podíváme jen na prostřední nejdelší sloupec:
\begin{center}
\includegraphics[width=0.8\textwidth]{generated/logic/multiplication/wallace-tree3.pdf}
\end{center}

\begin{itemize}
\item Vidíme, že za 8 kroků sčítaček, tedy 16 zpoždění hradel nám zbývá sečíst dva bity
\item Vpravo od tohoto sloupce je již něco sečteno a již se i propagovali carry posledních 8 sčítanců
\item Zbývá sečíst dvě 120-bitová čísla (součty a carry), což se dá stihnout asi za 30 zpoždění hradel
\item Výsledek - vynásobíme dvě čísla za cenu doby odpovídající přibližně dvěma součtům 64-bitových čísel
\item V praxi to trvá trochu déle, je použito méně sčítaček CSA.
\end{itemize}

\end{frame}



\begin{frame}
\frametitle{Testovací kvíz}

Kolik jste toho zatím pochopili?
\begin{itemize}
\item[A] Vše bez problému.
\item[B] Téměř vše.
\item[C] Téměř nic.
\item[D] Vůbec nic.
\end{itemize}

\end{frame}


\end{document}

